\subsection{独立な代数体例：中心部分環上階数81の完全な有限証明書}
\label{subsec:cubic-third-root-unit-boundary-L-finite-center}
\begin{proposition}[固定代数体例の中心次数と自由加群]
\label{prop:cubic-third-root-unit-boundary-L-finite-center}
\[
 K=\mathbb Q[t]/(t^6-9t^3+57),\quad
 d=t,\quad E=2,\quad\omega=(1-t^3)/7,\quad u=1,v=Z=0
\]
とする。この多項式は既約であり、全六つの埋込みで $L=0,\Omega\ne0$。
$C=B^\theta$、$A=C/(\widetilde\Gamma)$ と置く。
その中心次数は
\[
\begin{array}{c|ccccc}
 n&1&2&3&6&9\\\hline
 \dim Z(C)_n&1&1&2&3&7\\
 \dim Z(A)_n&1&1&2&3&5
\end{array}
\]
である。$C$ の $\delta^2$ 正常六次空間は二次元で、$\widetilde\Gamma^2$ 以外の元を持つ。
さらに、certificateに元の語の係数で明示した中心 $z_9\in A_9$ があり、
\[
 K[x,\Psi,z_9]\simeq K[X_1,X_3,X_9],\qquad
 A\text{ はこの部分環上階数81の自由加群}
\]
となる。これはこの代数体例の主張であり、全パラメータ又は全中心の同定ではない。
階数81だけからPI次数9や整域性を結論しない。
\end{proposition}
\begin{proof}
係数を最小多項式で厳密計算し、$\omega^2+\omega+1=0$ と $\Omega(2)\ne0$ を確認する。
$C$ の二次関係は $B$ の語 $ab$ の係数を
$\theta(b)$ の固有値で割って得る。
三次の語では $\omega^{\sum_{j=0}^2j\operatorname{wt}(a_j)}$ で割る。
この規約で元の六関係と $\widetilde\Gamma,\Psi$ を再生する。

高次数の全語 $4^n$ を展開する代わりに、既知の商基底と右乗法から
\[
 A_n=(A_{n-1}\otimes V)\Big/
 \sum_{r\in\{2,3\}}\operatorname{im}(A_{n-r}\otimes I_r)
\tag{*}
\]
を再帰的に計算する。
$C$ では $I_3$ を入れない。
接尾部で終わらない関係は既に $A_{n-1}$ で除かれているので、$(*)$ は全関係を使う厳密な商である。
基底語の左乗法も再帰的に回収する。
元の二次・三次関係の左右作用が、計算した全次数で零になることを別に確認する。

次数2から10の商次元は
\[
\begin{array}{c|rrrrrrrrr}
n&2&3&4&5&6&7&8&9&10\\\hline
C&10&20&35&56&84&120&165&220&286\\
A&10&19&31&46&64&85&109&136&166
\end{array}
\]
である。
各中心は全左右乗法行列の差の共通核で求める。
$C$ の中心は正則正常な $\widetilde\Gamma$ の正常化 $\delta$ で固定されるので、
その場合だけ固定文字部分に制限できる。
$A$ の中心はその制限を使わず全基底で計算する。
得た核基底は全交換子への再代入と独立性も確認し、表の次元を得る。
同様に正常六次の行列は $115\times30$、rank28である。

$Z(A)_9$ の中の四つの独立元
\[
 x^9,\quad x^6\Psi,\quad x^3\Psi^2,\quad\Psi^3
\]
に加え、五番目の独立元を $z_9$ として選ぶ。
元の語表示と正規化はcertificateに全て保持する。
$x,\Psi$ は $A$ 上で中心的正則列である。
実際、前節の通り $x$ は正則であり、$x=0$ で $\Psi$ はcanonical $g_0$ の倍。
$S_C/(g_0)$ はTHCR整域で、$\Gamma_0$ の像が非零なので、
正常正則列を交換すると $g_0$ は $S_C/(\Gamma_0)$ 上正則。
これを $x$ 進に持ち上げれば主張を得る。

$Q=A/(x,\Psi)$ と置くと
\[
 H_Q(t)=\frac{(1-t^3)^2}{(1-t)^3}.
\]
$z_9$ の $Q_9$ への像は非零。
さらに $Q/(z_9)$ を元の三生成元関係から次数13まで完全に消去すると、
次数0から13の次元は
\[
 1,3,6,8,9,9,9,9,9,8,6,3,1,0.
\]
一次生成なので次数13が零なら以後も零。
従って全Hilbert級数は正確に
\[
 H_{Q/(z_9)}(t)=(1+t+t^2)^2(1+t+\cdots+t^8)
              =(1-t^9)H_Q(t).
\]
$z_9$ 倍の完全列から、そのannihilatorのHilbert級数は零、すなわち $z_9$ は $Q$ 上正則。
これは次数13までの標本を全次数へ外挿したのではなく、一次生成による零次数以後の消滅を使った証明である。

中央の三元による商が81次元なので、次数帰納法から $A$ は
$S=K[x,\Psi,z_9]$ 上有限生成。
$\operatorname{GKdim} A=3$ と $S$ が三生成元であることから、それらは代数的独立。
商の81個の斉次基底を持ち上げて作る自由 $S$ 加群は $A$ へ全射する。
上のHilbert等式から両辺のHilbert級数は一致するので同型である。
従って自由階数は81。
完全な中心又は一般点の中心単純代数の次数は、なお別途判定する必要がある。
\end{proof}
